% Generated draft: requires mathematical and source review.
\documentclass[11pt]{article}
\usepackage[margin=25mm]{geometry}
\usepackage{fontspec}
\setmainfont{Latin Modern Roman}
\usepackage{amsmath,amssymb,mathtools,unicode-math}
\setmathfont{Latin Modern Math}
\usepackage{enumitem,needspace,xurl,hyperref}
\setlength{\parindent}{0pt}
\setlength{\parskip}{5pt}
\setlength{\emergencystretch}{3em}
\newcommand{\N}{\mathbb{N}}
\newcommand{\Z}{\mathbb{Z}}
\newcommand{\Q}{\mathbb{Q}}
\newcommand{\R}{\mathbb{R}}
\newcommand{\C}{\mathbb{C}}
\newcommand{\F}{\mathbb{F}}
\newcommand{\A}{\mathbb{A}}
\newcommand{\PP}{\mathbb{P}}
\newcommand{\E}{\mathbb{E}}
\newcommand{\eps}{\varepsilon}
\newcommand{\dd}{\,\mathrm{d}}
\begin{document}
\section*{UnsolvedMath ID: 11000116}
\textbf{Source problem code:} catalogue entry
\subsection{Definitions and mathematical statement}
Let $S$ be an oriented genus-$g$ surface with $g\geq3$ and $b\geq1$ boundary components. In the standard relative-boundary convention, $\operatorname{Mod}_g^b$ consists of orientation-preserving mapping classes fixing the boundary pointwise. Let $a_1,a_2,\ldots$ be any infinite sequence of nonseparating simple closed curves, let $t_a$ denote the right Dehn twist, and put $P_n=t_{a_1}t_{a_2}\cdots t_{a_n}$. Curves may repeat. The function $c$ occurring in the source is not defined in this record, so the following statements remain conditional on its preceding definition.

Problem 2.11 asks whether, for every such sequence,
\[
 \lim_{n\to\infty}\frac{c(P_n)}n
\]
exists and, if it exists, is always positive. It further asks for a positive lower bound $c(P_n)\geq nK$ for every $n$. The subsequent discussion singles out the stronger uniform question
\[
 \exists K=K(g,b)>0\quad\forall(a_i)_{i\geq1}\quad
 \forall n\geq1,\qquad c(P_n)\geq nK.
\]
The earlier lower-bound sentence does not explicitly say whether its $K$ may depend on the sequence. Existence of a positive limit alone also does not supply a stated all-$n$ bound without further information on $c$.

The source observes that such a sequence-independent $K$ would bound the length of any displayed positive nonseparating-twist factorization $f=t_{a_1}\cdots t_{a_k}$ by $k\leq c(f)/K$, and says this would answer the omitted Problem 2.9 positively. It emphasizes that $b\geq1$ is essential: on a closed surface a positive identity factorization can be repeated to give the reported zero limit, using the properties of $c$ implicit in the source. Conversely it reports no nonempty positive identity factorization for the boundary-fixed group, citing an embedding in a larger closed surface and Smith [41], whose theorem is described as saying that curves in a positive identity factorization fill the closed surface. The source credits Rostislav Matveyev, June 2000, for communicating that fact.
\subsection{Short English statement}
For arbitrary products of successive positive nonseparating Dehn twists on a surface with boundary, does the source's complexity divided by word length have a positive limit, and is there a uniform positive linear lower bound? The definition of the complexity is missing.
\subsection{Sources}
\begin{itemize}
\item Local registry: \url{data/unsolvedmath/problems.json}, permanent UnsolvedMath ID $11000116$, exact source code \texttt{catalogue entry}.
\item Farb, \emph{Problems on Mapping Class Groups and Related Topics}, Problem 2.11, PDF page 97: \url{https://www.math.uchicago.edu/~farb/papers/mcgbook.pdf}. Supplied access date: 2026-07-29; extraction: pdf-font-structure; repeated local numbering disambiguated by page among 329 headings.
\item Supplied citation: Smith [41], with the filling statement credited as learned from Rostislav Matveyev in June 2000. Problem 2.9 and the definition of $c$ are absent, as is the full bibliography entry. The registry literature summary is truncated.
\end{itemize}
\end{document}
